\section*{Capítulo \ref{ch:g8:fronts-and-weather-maps} --- Masas de aire, frentes y mapas del tiempo}

\begin{solution}{exo:g8:fronts-and-weather-maps:1}
(a) Tropical marítima: cálida y húmeda. (b) Polar continental: fría y
seca. (c) Tropical continental: muy cálida y seca. (d) Polar marítima:
fría y húmeda.
\end{solution}

\begin{solution}{exo:g8:fronts-and-weather-maps:2}
Línea azul con triángulos: un \omterm{def:g8:fronts-and-weather-maps:front}{frente frío}. Línea roja con semicírculos: un
\omterm{def:g8:fronts-and-weather-maps:front}{frente cálido}. Línea que une lugares con la misma presión: una \omterm{def:g8:fronts-and-weather-maps:isobar}{isobara}. A:
el centro de un \omterm{def:g8:fronts-and-weather-maps:depression}{anticiclón} (altas presiones). B: el centro de una \omterm{def:g8:fronts-and-weather-maps:depression}{borrasca}
(bajas presiones).
\end{solution}

\begin{solution}{exo:g8:fronts-and-weather-maps:3}
(a) Verdadero. (b) Falso: llegan antes que él, mucho antes, muy altas en
el cielo. (c) Falso: las \omterm{def:g8:fronts-and-weather-maps:isobar}{isobaras} juntas indican un viento fuerte.
(d) Verdadero.
\end{solution}

\begin{solution}{exo:g8:fronts-and-weather-maps:4}
(a) \qty{980}{hPa}. (b) \qty{1034}{hPa}. (c) \qty{4}{hPa}. (d) $1034 -
980 = \qty{54}{hPa}$.
\end{solution}

\begin{solution}{exo:g8:fronts-and-weather-maps:5}
Finas estrías de \omterm{def:g4:clouds-rain-snow:cloud-families}{cirros}; un velo gris de altostratos; las \omterm{def:g4:clouds-rain-snow:cloud}{nubes} bajan y se
espesan; \omterm{def:g4:clouds-rain-snow:precipitation}{lluvia} continua; aire más cálido y más húmedo.
\end{solution}

\begin{solution}{exo:g8:fronts-and-weather-maps:6}
En un \omterm{def:g8:fronts-and-weather-maps:front}{frente frío}, el aire frío, denso, se mete bajo el aire cálido como
una \emph{cuña}: el aire cálido se eleva por una pendiente
\emph{empinada}, sube deprisa y forma \omterm{def:g4:clouds-rain-snow:cloud}{nubes} altas. En un \omterm{def:g8:fronts-and-weather-maps:front}{frente cálido}, el
aire cálido sube lentamente por encima del aire frío por una
\emph{pendiente suave}: se eleva poco a poco a lo largo de una gran
distancia, y sus \omterm{def:g4:clouds-rain-snow:cloud}{nubes} se extienden en capas.
\end{solution}

\begin{solution}{exo:g8:fronts-and-weather-maps:7}
(a) Cerca de B: allí las \omterm{def:g8:fronts-and-weather-maps:isobar}{isobaras} están más juntas y las flechas son más
largas. (b) Del este (las flechas apuntan al oeste, un poco hacia el
sudoeste). (c) Del este (las flechas apuntan al oeste). (d) En sentido
contrario a las agujas del reloj.
\end{solution}

\begin{solution}{exo:g8:fronts-and-weather-maps:8}
(a) $1016 - 1003 = \qty{13}{hPa}$. (b) Una \omterm{def:g8:fronts-and-weather-maps:depression}{borrasca}: la presión no deja de
bajar. (c) \omterm{def:g4:clouds-rain-snow:cloud}{Nubes}, viento y \omterm{def:g4:clouds-rain-snow:precipitation}{lluvia}, probablemente con un \omterm{def:g8:fronts-and-weather-maps:front}{frente}.
\end{solution}

\begin{solution}{exo:g8:fronts-and-weather-maps:9}
(a) cT, tropical continental. (b) Muy cálida y seca. (c) Sobre el \omterm{def:g2:water-on-earth:ocean}{mar}
recoge vapor de agua: se vuelve más húmeda, y el aire cercano al agua se
enfría un poco al contacto con el \omterm{def:g2:water-on-earth:ocean}{mar}.
\end{solution}

\begin{solution}{exo:g8:fronts-and-weather-maps:10}
(a) $4 - 1 = 3$ días. (b) $3 \div 3 = 1$ día por cada diez años. (c) Hacen
falta $5 - 1 = 4$ días más: unos 40 años.
\end{solution}

\begin{solution}{exo:g8:fronts-and-weather-maps:11}
(a) Unos \qty{1350}{km} (cuatro barras de \omterm{def:g2:drawing-to-map:scale}{escala} y media).
(b) \qty{4}{hPa} en \qty{80}{km}: $4 \div 80 \times 100 = \qty{5}{hPa}$ por
\qty{100}{km}. (c) $4 \div 140 \times 100 \approx \qty{2.9}{hPa}$ por
\qty{100}{km}. (d) Cerca de B, donde la presión cambia más deprisa; sí,
las flechas son más largas cerca de B.
\end{solution}

\begin{solution}{exo:g8:fronts-and-weather-maps:12}
(a) El \omterm{def:g8:fronts-and-weather-maps:front}{frente frío}: $300 \div 50 = 6$ horas. El \omterm{def:g8:fronts-and-weather-maps:front}{frente cálido} avanza a
\qty{25}{km} por hora: $200 \div 25 = 8$ horas. El \omterm{def:g8:fronts-and-weather-maps:front}{frente frío} llega
2 horas antes. (b) Primera ciudad: \omterm{def:g4:clouds-rain-snow:cloud}{nubes} de gran desarrollo vertical,
chubascos fuertes, rachas, y después aire más fresco y un cielo que se
despeja. Segunda ciudad: \omterm{def:g4:clouds-rain-snow:cloud}{nubes} en capas cada vez más bajas y una \omterm{def:g4:clouds-rain-snow:precipitation}{lluvia}
larga y continua, y después aire más cálido y más húmedo. (c) La segunda
ciudad, a la que llega el aire cálido.
\end{solution}

\begin{solution}{pb:g8:fronts-and-weather-maps:1}
\textbf{1.} En el aire cálido, entre los dos \omterm{def:g8:fronts-and-weather-maps:front}{frentes}. El \omterm{def:g8:fronts-and-weather-maps:front}{frente cálido} ya
ha pasado.

\textbf{2.} Del sudoeste.

\textbf{3.} Hacia el este: los triángulos de un \omterm{def:g8:fronts-and-weather-maps:front}{frente frío} están del lado
hacia el que avanza, y apuntan al este, hacia la ciudad.

\textbf{4.} \omterm{def:g4:clouds-rain-snow:cloud}{Nubes} de gran desarrollo vertical, chubascos fuertes, quizá
truenos y \omterm{def:g4:clouds-rain-snow:precipitation}{granizo}, rachas de viento; después, aire más fresco, viento del
oeste o del noroeste y un cielo que se despeja.

\textbf{5.} A \qty{180}{km} a las 06:00 y a \qty{300}{km} a las 09:00.

\textbf{6.} $300 - 180 = \qty{120}{km}$.

\textbf{7.} $120 \div 3 = \qty{40}{km}$ por hora.

\textbf{8.} La ciudad está a \qty{480}{km}: $480 - 300 = \qty{180}{km}$.

\textbf{9.} $180 \div 40 = 4.5$ horas, es decir, 4 horas y 30 minutos.

\textbf{10.} 09:00 + 4~h~30~min = \textbf{13:30}.

\textbf{11.} No: los chubascos deberían terminar hacia las 14:30, antes del
comienzo del partido.

\textbf{12.} $180 \div 30 = 6$ horas: a las 15:00, justo al comienzo del
partido. Previsión: ``El \omterm{def:g8:fronts-and-weather-maps:front}{frente frío} debería llegar a la ciudad hacia las
\textbf{13:30}, con una hora de chubascos fuertes y rachas, y después
cielos despejados: lo más probable es que el partido pueda jugarse. Si el
\omterm{def:g8:fronts-and-weather-maps:front}{frente} frena, podría llegar a las 15:00 y retrasar el partido''.
\end{solution}
